\documentclass{article}

\usepackage[utf8]{inputenc}
\usepackage[spanish,es-noshorthands]{babel}
\usepackage{amssymb, amsmath, amsbsy}
\usepackage{tasks}
\usepackage{graphicx}

\fontfamily{cmss}\selectfont
\renewcommand{\familydefault}{\sfdefault}
\pagestyle{empty}

\begin{document}

\section*{Relaciones binarias}

\noindent Sean $A$ y $B$ dos conjuntos no vacíos. Un conjunto $\Re$ de pares ordenados, se llama relación de A en B, si $\Re$ es un subconjunto cualquiera de $A\times B$, es decir : \\

$\Re$ es una relación de $A$ en $B$, si y sólo si $\Re \subset A\times B$.

% \paragraph*{Ejemplo} -- Sean
\end{document}